%\clearpage
\section{Background on Information Theory}\label{prelim-sec:info}

We now briefly introduce some definitions and facts from information theory that are needed in this thesis. We refer the interested reader to the text by Cover and Thomas~\cite{CoverT06} for an excellent introduction to this field, 
and the proofs of the statements used in this Appendix. 

For a random variable $\rA$, we use $\supp{\rA}$ to denote the support of $\rA$ and $\distribution{\rA}$ to denote its distribution. 
When it is clear from the context, we may abuse the notation and use $\rA$ directly instead of $\distribution{\rA}$, for example, write 
$A \sim \rA$ to mean $A \sim \distribution{\rA}$, i.e., $A$ is sampled from the distribution of random variable $\rA$. 


We denote the \emph{Shannon Entropy} of a random variable $\rA$ by
$\en{\rA}$, which is defined as: 
\begin{align}
	\en{\rA} := \sum_{A \in \supp{\rA}} \Pr\paren{\rA = A} \cdot \log{\paren{1/\Pr\paren{\rA = A}}} \label{eq:entropy}
\end{align} 

\noindent
The \emph{conditional entropy} of $\rA$ conditioned on $\rB$ is denoted by $\en{\rA \mid \rB}$ and defined as:
\begin{align}
\en{\rA \mid \rB} := \E_{B \sim \rB} \bracket{\en{\rA \mid \rB = B}}, \label{eq:cond-entropy}
\end{align}
where 
$\en{\rA \mid \rB = B}$ is defined in a standard way by using the distribution of $\rA$ conditioned on the event $\rB = B$ in Eq~(\ref{eq:entropy}). 

\noindent
Recall that the binary entropy $H(p)$ is the entropy of a two-valued random
variable that takes one value with probability $p$.
We need the following upper bound on the binary entropy.

\begin{fact}\label{fact:ent-ub}
  The binary entropy $H_2(p) \leq 2 \sqrt{ p \cdot (1 - p) }$.
\end{fact}


 
The \emph{mutual information} of two random variables $\rA$ and $\rB$ is denoted by
$\mi{\rA}{\rB}$ and is defined as:
\begin{align}
\mi{\rA}{\rB} := \en{A} - \en{A \mid  B} = \en{B} - \en{B \mid  A}. \label{eq:mi}
\end{align}
\noindent
The \emph{conditional mutual information} $\mi{\rA}{\rB \mid \rC}$ is $\en{\rA \mid \rC} - \en{\rA \mid \rB,\rC}$ and hence by linearity of expectation:
\begin{align}
	\mi{\rA}{\rB \mid \rC} = \E_{C \sim \rC} \bracket{\mi{\rA}{\rB \mid \rC = C}}. \label{eq:cond-mi}
\end{align} 



\subsection{Useful Properties of Entropy and Mutual Information}\label{sec:prop-en-mi}

We shall use the following basic properties of entropy and mutual information throughout. 


\begin{fact}\label{fact:it-facts}
  Let $\rA$, $\rB$, $\rC$, and $\rD$ be four (possibly correlated) random variables.
   \begin{enumerate}
  \item \label{part:uniform} $0 \leq \en{\rA} \leq \log|\supp{\rA}|$. The right equality holds
    iff $\distribution{\rA}$ is uniform.
  \item \label{part:info-zero} $\mi{\rA}{\rB \mid \rC} \geq 0$. The equality holds iff $\rA$ and
    $\rB$ are \emph{independent} conditioned on $\rC$.
    \item \label{part:info-atmost-rv} $\mi{\rA}{\rB \mid \rC} \leq \en{\rA}$.
  \item \label{part:cond-reduce} \emph{Conditioning on a random variable reduces entropy}:
    $\en{\rA \mid \rB,\rC} \leq \en{\rA \mid  \rB}$.  The equality holds iff $\rA \perp \rC \mid \rB$.
    \item \label{part:sub-additivity} \emph{Subadditivity of entropy}: $\en{\rA,\rB \mid \rC}
    \leq \en{\rA \mid C} + \en{\rB \mid  \rC}$.
   \item \label{part:ent-chain-rule} \emph{Chain rule for entropy}: $\en{\rA,\rB \mid \rC} = \en{\rA \mid \rC} + \en{\rB \mid \rC,\rA}$.
  \item \label{part:chain-rule} \emph{Chain rule for mutual information}: $\mi{\rA,\rB}{\rC \mid \rD} = \mi{\rA}{\rC \mid \rD} + \mi{\rB}{\rC \mid  \rA,\rD}$.
  \item \label{part:data-processing} \emph{Data processing inequality}: for a function $f(\rA)$ of $\rA$, $\mi{f(\rA)}{\rB \mid \rC} \leq \mi{\rA}{\rB \mid \rC}$. 
   \end{enumerate}
\end{fact}


\noindent
We also use the following two standard propositions, regarding the effect of conditioning on mutual information.

\begin{proposition}\label{prop:info-increase}
  For random variables $\rA, \rB, \rC, \rD$, if $\rA \perp \rD \mid \rC$, then, 
  \[\mi{\rA}{\rB \mid \rC} \leq \mi{\rA}{\rB \mid  \rC,  \rD}.\]
\end{proposition}
 \begin{proof}
  Since $\rA$ and $\rD$ are independent conditioned on $\rC$, by
  \itfacts{cond-reduce}, $\HH(\rA \mid  \rC) = \HH(\rA \mid \rC, \rD)$ and $\HH(\rA \mid  \rC, \rB) \ge \HH(\rA \mid  \rC, \rB, \rD)$.  We have,
	 \begin{align*}
	  \mi{\rA}{\rB \mid  \rC} &= \HH(\rA \mid \rC) - \HH(\rA \mid \rC, \rB) = \HH(\rA \mid  \rC, \rD) - \HH(\rA \mid \rC, \rB) \\
	  &\leq \HH(\rA \mid \rC, \rD) - \HH(\rA \mid \rC, \rB, \rD) = \mi{\rA}{\rB \mid \rC, \rD}. \qed
	\end{align*}
	
\end{proof}

\begin{proposition}\label{prop:info-decrease}
  For random variables $\rA, \rB, \rC,\rD$, if $ \rA \perp \rD \mid \rB,\rC$, then, 
  \[\mi{\rA}{\rB \mid \rC} \geq \mi{\rA}{\rB \mid \rC, \rD}.\]
\end{proposition}
 \begin{proof}
 Since $\rA \perp \rD \mid \rB,\rC$, by \itfacts{cond-reduce}, $\HH(\rA \mid \rB,\rC) = \HH(\rA \mid \rB,\rC,\rD)$. Moreover, since conditioning can only reduce the entropy (again by \itfacts{cond-reduce}), 
  \begin{align*}
 	\mi{\rA}{\rB \mid  \rC} &= \HH(\rA \mid \rC) - \HH(\rA \mid \rB,\rC) \geq \HH(\rA \mid \rD,\rC) - \HH(\rA \mid \rB,\rC) \\
	&= \HH(\rA \mid \rD,\rC) - \HH(\rA \mid \rB,\rC,\rD) = \mi{\rA}{\rB \mid \rC,\rD}. \qed
 \end{align*}

\end{proof}


\subsection{Background on Communication and Information Complexity}\label{sec:cc-ic}

\paragraph{Communication complexity.} We briefly review the standard definitions of the two-party communication model of Yao~\cite{Yao79}.  See the text by Kushilevitz and Nisan~\cite{KushilevitzN97} for 
an extensive overview of communication complexity. In Section~\ref{sec:lower-bound-hpc}, we also use a standard generalization of this model to allow for more than two players, but we defer
the necessary definitions to that section. 

Let $P: \mathcal{X} \times \mathcal{Y} \rightarrow \mathcal{Z}$ be a relation.  Alice receives an input $X
\in \mathcal{X}$ and Bob receives $Y \in \mathcal{Y}$, where $(X,Y)$ are chosen from a
joint distribution $\cD$ over $\mathcal{X} \times \mathcal{Y}$. We allow players to have access to both public and private randomness. 
They communicate with each other by exchanging messages such that each message
depends only on the private input and random bits of the player sending the message, and the already communicated messages plus the public randomness. 
At the end, one of the players need to output an answer  $Z$ such that $Z \in P(X,Y)$.  

We use $\prot$ to denote a protocol used by the players. We always assume that the protocol $\prot$ can be randomized (using both public and
private randomness), \emph{even against a prior distribution $\cD$ of inputs}. For any
$0 < \delta < 1$, we say $\prot$ is a $\delta$-error protocol for $P$ over a distribution
$\cD$, if the probability that for an input $(X,Y)$, $\prot$ outputs some $Z$ where $Z \notin P(X,Y)$ is at most
$\delta$ (the probability is taken over the randomness of \emph{both} the distribution and the protocol).

\begin{definition}[Communication cost]
  The \emph{communication cost} of a protocol $\prot$ on an input
  distribution $\cD$, denoted by $\CC{\prot}{\cD}$, is the \emph{worst-case} bit-length of the transcript
  communicated between Alice and Bob in the protocol $\prot$, when the inputs are chosen from $\cD$.

\end{definition}

Communication complexity of a problem $P$ is defined as the minimum communication cost of a protocol $\prot$ that solves $P$ on every distribution $\cD$ with probability at least $2/3$. 



\paragraph{Information complexity.} There are several possible definitions of information
cost of a communication prtocol that have been considered depending on the application (see, e.g.,~\cite{ChakrabartiSWY01,Bar-YossefJKS02,BarakBCR10,BravermanR11,BravermanEOPV13}).  
We use the notion of \emph{internal information cost}~\cite{BarakBCR10} that measures the average amount of information each player learns about the input of the other player
by observing the transcript of the protocol. %Formally, 
\begin{definition}[Information cost]
  Consider an input distribution $\cD$ and a protocol $\prot$. Let $(\rX,\rY) \sim \cD$ denote the random variables for the input of Alice and Bob and $\rProt$ be the 
  the random variable for the transcript of the protocol \emph{concatenated} with the public randomness $\rR$ used by $\prot$. 
  The \emph{(internal) information cost} of $\prot$ with respect to
  $\cD$ is $\ICost{\prot}{\cD}:=\mii{\rProt}{\rX \mid \rY}{\cD} + \mii{\rProt}{\rY \mid \rX}{\cD}$. 

\end{definition}

One can also define information complexity of a problem $P$ similar to communication complexity with respect to the information cost. However, we avoid presenting this definition formally due to some subtle technical issues
 that need to be addressed which lead to multiple different but similar-in-spirit definitions. As such, we state our results directly in terms of information cost. 

%%
Note that any public coin protocol is a distribution over private coins protocols, run by
first using public randomness to sample a random string $\rR=R$ and then running the
corresponding private coin protocol $\prot^R$. We also use $\rProt^R$ to denote the transcript of the protocol $\prot^R$. 
We have the following standard proposition. 
\begin{proposition}\label{prop:public-random}
	 For any distribution $\cD$ and any protocol $\prot$ with public randomness $\bR$, $$\ICost{\prot}{\cD} = \mii{\rProt}{\rX \mid \rY, \rR}{\cD} + \mii{\rProt}{\rY \mid \rX, \rR}{\cD} = \E_{R \sim \rR}\bracket{\ICost{\prot^R}{\cD}}.$$ 
\end{proposition}
\begin{proof} By definition of internal information cost, 
\begin{align*}
	\ICost{\prot}{\cD} &= \mii{\rProt}{\rX \mid \rY}{\cD} + \mii{\rProt }{\rY \mid \rX}{\cD} = \mi{\rProt,\rR}{\rX \mid \rY} + \mi{\rProt,\rR}{\rY \mid \rX} \tag{$\Prot$ denotes the transcript and the public randomness} \\
	&= \mi{\rR}{\rX \mid \rY}  + \mi{\rProt}{\rX \mid \rY , \rR} + \mi{\rR}{\rY \mid \rX}  + \mi{\rProt}{\rY \mid \rX , \rR}  \tag{chain rule of mutual information, \itfacts{chain-rule}} \\
	&= \mi{\rProt}{\rX \mid \rY , \rR} + \mi{\rProt}{\rY \mid \rX , \rR} \tag{$\mi{\rR}{\rX \mid \rY}  = \mi{\rR}{\rY \mid \rX}  = 0$ since $\rR \perp \rX, \rY$ and \itfacts{info-zero}} \\
	&= \E_{R \sim \rR}\bracket{\mi{\rProt}{\rX \mid \rY , \rR=R} + \mi{\rProt}{\rY \mid \rX , \rR=R}} = \E_{R \sim \rR}\bracket{\ICost{\prot^R}{\cD}}, 
\end{align*}
concluding the proof.
\end{proof}


The following well-known proposition relates communication cost and information cost. 
\begin{proposition}[cf.~\cite{BravermanR11}]\label{prop:cc-ic}
  For any distribution $\cD$ and any protocol $\prot$: $\ICost{\prot}{\cD} \leq \CC{\prot}{\cD}$.
\end{proposition}
\begin{proof}
Let us assume first that $\prot$ only uses private randomness and thus $\rProt$ only contain the transcript.  
For any $b \in [\CC{\prot}{\cD}]$, we define $\Prot_b$ to be the $b$-th bit of the transcript. We have, 
\begin{align*}
	 \ICost{\prot}{\cD} &= \mi{\rProt}{\rX \mid \rY}+ \mi{\rProt }{\rY \mid \rX} \\
	 &= \sum_{b=1}^{\CC{\prot}{\cD}} \mi{\rProt_b}{\rX \mid \rProt^{<b},\rY} + \mi{\rProt_b}{\rY \mid \rProt^{<b},\rX} \tag{by chain rule of mutual information in \itfacts{chain-rule}}\\
	 &= \sum_{b=1}^{\CC{\prot}{\cD}} \E_{\Prot^{<b}}\bracket{\mi{\rProt_b}{\rX \mid \rProt^{<b} = \Prot^{<b},\rY} + \mi{\rProt_b}{\rY \mid \rProt^{<b}= \Prot^{<b},\rX}}.
\end{align*}
Consider each term in the RHS above. By conditioning on $\Prot^{<b}$, the player that transmit $\rProt_b$ would become fix. If this player is Alice, then $\mi{\rProt_b}{\rY \mid \rProt^{<b}= \Prot^{<b},\rX} = 0$, because
$\rProt_b$ is only a function of $(\rProt^{<b},\rX)$ in this case; similarly, if this player is Bob, then $\mi{\rProt_b}{\rX \mid \rProt^{<b}= \Prot^{<b},\rY} = 0$. Moreover, 
$\mi{\rProt_b}{\rX \mid \rProt^{<b} = \Prot^{<b},\rY} \leq \en{\rProt_b}  \leq 1$ and similarly $\mi{\rProt_b}{\rY \mid \rProt^{<b}= \Prot^{<b},\rX} \leq 1$. As such, the above term can be upper bounded 
by $\CC{\prot}{\cD}$. To finalize the proof, note that by Proposition~\ref{prop:public-random}, for any public-coin protocol $\prot$, $\ICost{\prot}{\cD} = \E_{R \sim \rR}\bracket{\ICost{\prot^R}{\cD}} \leq \E_{R \sim \rR}\bracket{\CC{\prot^R}{\cD}} \leq \CC{\prot}{\cD}$, where
the first inequality is by the first part of the argument.  
\end{proof}

Proposition~\ref{prop:cc-ic} provides a convinent way of proving communication complexity lower bounds by lower bounding information cost of any protocol. 


\subsection*{Rectangle Property of Communication Protocols}\label{sec:statstics}
We conclude this section by mentioning some basic properties of communication protocols. For any protocol $\prot$ and inputs $x \in \mathcal{X}$ and $y \in \mathcal{Y}$, we define $\Prot_{x,y}$ as the transcript
of the protocol conditioned on the input $x$ to Alice and input $y$ to Bob. Note that for randomized protocols, $\Prot_{x,y}$ is a random variable which we denote by $\rProt_{x,y}$. 

The following is referred to as the rectangle property of deterministic protocols. 
\begin{fact}[Rectangle property]\label{fact:rectangle}
	For any deterministic protocol $\prot$ and inputs $x,x' \in \mathcal{X}$ to Alice and $y,y' \in \mathcal{Y}$ to Bob, if $\Prot_{x,y} = \Prot_{x',y'}$, then $\Prot_{x,y'} = \Prot_{x',y}$. 
\end{fact}
\noindent
Fact~\ref{fact:rectangle} implies that the set of inputs consistent with any transcript $\Prot_{x,y}$ of a deterministic protocol forms a combinatorial rectangle. 
One can also extend the rectangle property of deterministic protocols to randomized protocols using the following fact. 
\begin{fact}[Cut-and-paste property; cf.~\cite{Bar-YossefJKS02}]\label{fact:r-rectangle}
	For any randomized protocol $\prot$ and inputs $x,x' \in \mathcal{X}$ to Alice and $y,y' \in \mathcal{Y}$ to Bob, $\hd{\rProt_{x,y}}{\rProt_{x',y'}} = \hd{\rProt_{x,y'}}{\rProt_{x',y}}$. 
\end{fact}
